\chapter{Prólogo}

\setlength{\parindent}{0cm} La decisión de haber entrado a estudiar Física en la Facultad de Ciencias de la UNAM ha sido de las más importantes de mi vida. Tuve una ardua trayectoria y una constante pregunta, sobre el sentido que le daría a mi formación académica. Aunque exploré diversas áreas como Cuántica y Relatividad principalmente, mis intereses se fueron orillando hacia la computación científica. Los cursos de métodos numéricos me abrieron camino para cursar otras materias que fueran más aplicadas, tal y como la que impartió el Dr. Sergio Alcalá ---posteriormente mi asesor de tesis--- sobre sistemas complejos.\\
\\
Este campo me permitió explorar una perspectiva interdisciplinaria de la ciencia, en donde la forma de abordar los fenómenos es partir de diversas perspectivas; particularmente nosotros como físicos desde el ámbito de la dinámica no lineal y la física estadística. Encontrar esta área significó mucho para mi a tal grado de buscar especializarme más en otros cursos como el de redes complejas con el Dr. Franciso Betancourt, el de dinámica no lineal de la Dra. Laura Palomares y el de introducción a sistemas complejos del Dr. Gustavo Martínez.\\
\\
Esta formación me permitió armar un criterio sobre mi tema de tesis ---el cual es sobre la estabilidad en el modelo poblacional de Lotka--Volterra--- y junto con la asesoría y guía del Dr. Sergio pudimos arrancar y desarrollar la investigación. Al principio el trabajo se centró en la parte que más me apasiona: desarrollar los algoritmos para poder desentrañar la dinámica del modelo generalizándolo a $N$ ecuaciones diferenciales. La pregunta de investigación desde un principio fue \textit{encontrar condiciones en las interacciones del modelo para garantizar estabilidad}.\\
\\
He de admitir que no fue una pregunta trivial, y considero que esta tesis solamente la resuelve parcialmente. En el trayecto fuimos descubriendo y aprendiendo nuevos aspectos de los sistemas ecológicos que normalmente no se enseña en un salón de clases. El trabajo de Robert May, Stefano Allesina, Guy Bunin, Giulio Biroli y Ada Altieri fueron fundamentales guías para este aprendizaje. Una vez que el desarrollo de los algoritmos satisficieron resultados óptimos, tras haberlos evaluado y sometido a constantes pruebas, tuvimos la oportunidad de correr simulaciones en un cluster ya que los tiempos de compilación fueron extenuantes.
\newline
Al obtener nuestros primeros resultados y redactar lo que fue el desarrollo de esta tesis, fuimos hallando nuevos descubrimientos sobre el comportamiento del sistema, lo que nos llevó a numerosas discusiones y revisión de nueva literatura para la interpretación de nuestras observaciones. En ese aspecto, pude encarar lo que significa el trabajo de investigación, que en analogía lo interpreto como la navegación en aguas misteriosas y su exploración constante para lograr entender un panorama cada vez más amplio del área. Por lo tanto ---y en honor a la no linealidad--- la redacción de esta tesis tomó numerosas revisiones, correcciones y replanteamientos, en esencia ha sido reescrita varias veces.\\
\\
Sin embargo, el trabajo logró su maduración cuando hallamos las justificaciones que sostienen los argumentos principales de esta tesis. Principalmente la \textit{factibilidad} en sistemas ecológicos fue el elemento que congregó todo y permitió la elaboración del argumento central que habla sobre un parámetro heurístico que delimita el umbral de estabilidad en el modelo LVG. Aún así, este trabajo requiere de esfuerzos mayores para poder hallar un criterio formal y espectral que delimite el umbral de estabilidad del sistema.\\
\\
Haber hallado amor y pasión sobre los sistemas complejos ha sido, en gran parte, debido a la formación y educación de mis padres. Ambos dedicados a las ciencias sociales, me heredaron esa chispa que se encuentra viva en la interdisciplinariedad que caracteriza a los sistemas complejos. En esta ocasión me ha tocado investigar y trabajar apasionadamente en ecología, y en algún futuro me proyecto a seguir estas investigaciones en el posgrado, ya sea sobre esta línea u otras afines.

\chapter{Agradecimientos}

Esta tesis fue desarrollada únicamente gracias al apoyo físico, mental y espiritual de los siguientes a mencionar:\\
\\
Principalmente agradezco a mi universidad y a mi facultad por brindarnos la formación y el cúmulo de caminos posibles por recorrer. Agradecer ampliamente a mi pueblo por sostener económicamente mi formación y la de todxs lxs estudiantes de esta universidad pública y gratuita.\\
\\
Agradezco ampliamente a mi asesor, el Dr. Sergio A. Alcalá Corona por haberme guiado acertadamente en todo el proceso de investigación, ha sido un pilar muy importante en mi formación académica. Agradezco a mi sínodo compuesto por los Dres. José Roberto Romero Arias, Octavio Raymundo Miramontes Vidal, Gustavo Carlos Martínez Mekler y Edgar Javier González Liceaga, por haber sido muy atentos y comprometidos con nuestro trabajo. Agradezco con cariño al Dr. Carlos Francisco Betancourt Moreno por haber contribuido de forma importante a mi formación y a la motivación de mi tesis. Asimismo agradezco la formación que obtuve de los cursos de los Dres. Octavio Miramontes y Gustavo Martínez sobre sistemas complejos.\\
\\
Agradezco a mis amigos Isaac y Marlene por haberme acompañado gran parte de la carrera, así como a Aline, Alexis, Diana, Betty y Osiris. Agradezco a Enya por acompañarme en mis últimos semestres de la carrera y parte del periodo laboral en Softtek.\\
\\
Agradezco a mi abuela Mayi por ser muy insistente en mi titulación pero también por todo el soporte y amor que me ha brindado siempre. A mi abuelo Efrén por el cariño y la atención de estar al pendiente de mi bienestar. A mi abuelo Armando por el amor y cuidado que me tiene y por todas esas veces que me llevó al metro Pantitlán. A Paty por todo el amor y atención que me tuvo. A mi abuela Carmen por tenerme presente a pesar de la distancia.
\newpage
Agradezco a mi tía Ara por todo el amor, cuidados y atenciones a mi salud, principalmente a mi audición. A mi tía Ixchel por ser egresada de la Facultad de Ciencias y por toda la guía que me dio. A mi tía Vinni por ser ejemplar en el área científica. A mi tío Alejandro, mi tía Yoli y mis primas Alexa y Fernanda, por ser tan calurosxs y cariñosxs conmigo. A mi tía Berenice, mi tío Eliseo y mi primo Santiago por ser tan atentxs y fraternxs conmigo. A mi tío Yax, mi tía Eli y mi primo Iker por ser tan alegres y fraternxs conmigo. A mis padrinxs Janet y Giovanni por ser tan bohemixs.\\
\\
A agradezco mi mamá Margarita por estar siempre presente, tanto físicamente como espiritualmente, por guiarme en las adversidades más grandes que he afrontado, por las incontables conversaciones que han elevado nuestro pensamiento, por todos los consejos y formación que contribuye en gran medida a la persona que soy hoy día. A mi papá Iván por heredarme el espíritu de lucha revolucionaria, la consciencia de clase y el amor por la movilidad y la salud física, por siempre echarme porras y alentar todo lo que hago. A mi hermano Fernando, por ser un excelente hermano y por todo lo que hemos vivido juntos desde el primer día que te conocí.\\
\\
Agradezco a Låttice, nuestro proyecto musical formado en 2025 integrado por Mora, Marco, Chema y Andrés, por toda esa música que hemos construido juntos y que ha servido de catarsis para la elaboración de esta tesis. A Mim por haberme acompañado desde 2018 y por ser testigo de todo el proceso de elaboración de esta tesis.\\
\\
Agradezco a Pao, por alentarme y prevenirme de aspectos clave del trámite de titulación. Y de esta forma, agradezco a todxs aquellxs que me guardan cariño y/o aprecio. El resultado de esta tesis no es únicamente un esfuerzo individual, es el producto del esfuerzo colectivo del entramado complejo al cual llamamos comunidad. Con ustedes todo, sin ustedes nada.\\
\\
\begin{flushright}%
\emph{Gracias.}
\end{flushright}